\documentclass{article}
\usepackage{amsmath,amssymb,amsthm}
\usepackage{algorithm}
\usepackage{algpseudocode}
\usepackage{enumitem}
\usepackage{hyperref}
\newtheorem{theorem}{Theorem}
\newtheorem{lemma}{Lemma}
\newtheorem{assumption}{Assumption}
\newtheorem{definition}{Definition}
\newcommand{\E}{\mathbb{E}}
\newcommand{\R}{\mathbb{R}}

\begin{document}

\section*{Probabilistic Momentum Coordinate Descent (PMCD)}

\section*{Mathematical Formulation}
Let $f:\R^{d}\to\R$ be continuously differentiable.  
For each iteration $k\ge 0$:

\begin{enumerate}[label=\arabic*)]
\item Sample a coordinate $i_k\in\{1,\dots ,d\}$ according to a fixed probability
      distribution $p=(p_1,\dots ,p_d)$ with $p_i>0$ and $\sum_{i=1}^{d}p_i=1$.
\item Compute the partial derivative (coordinate gradient)
      \[
          g_{i_k}^{(k)}:=\nabla_{i_k}f(w^{(k)}) .
      \]
\item Update the momentum term for the selected coordinate:
      \[
          v_{i_k}^{(k+1)} \;=\; \beta\, v_{i_k}^{(k)} \;+\; (1-\beta)\, g_{i_k}^{(k)},
      \qquad \beta\in[0,1).
      \tag{1}
      \]
\item Perform a coordinatewise gradient step using the momentum:
      \[
          w_{i_k}^{(k+1)} \;=\; w_{i_k}^{(k)} \;-\; \eta\, v_{i_k}^{(k+1)},
      \qquad \eta>0,
      \tag{2}
      \]
      while all other coordinates remain unchanged,
      $w_{j}^{(k+1)}=w_{j}^{(k)}$ for $j\neq i_k$.
\end{enumerate}

All vectors $w^{(k)}\in\R^{d}$ are initialized arbitrarily, and the momentum
vector $v^{(0)}$ is usually set to $\mathbf 0$.

\section*{Pseudocode}
\begin{algorithm}[H]
\caption{Probabilistic Momentum Coordinate Descent (PMCD)}
\begin{algorithmic}[1]
\State \textbf{Input:} step size $\eta>0$, momentum $\beta\in[0,1)$, probability vector $p\in\Delta^{d-1}$, max. iter $T$.
\State Initialize $w^{(0)}\in\R^{d}$, $v^{(0)}\gets\mathbf 0$.
\For{$k=0$ \textbf{to} $T-1$}
    \State Sample $i_k\sim p$.
    \State $g_{i_k}^{(k)}\gets \nabla_{i_k} f(w^{(k)})$.
    \State $v_{i_k}^{(k+1)} \gets \beta\,v_{i_k}^{(k)}+(1-\beta)g_{i_k}^{(k)}$ \Comment{momentum update (1)}
    \State $w_{i_k}^{(k+1)} \gets w_{i_k}^{(k)}-\eta\,v_{i_k}^{(k+1)}$ \Comment{coordinate step (2)}
    \State $w_{j}^{(k+1)}\gets w_{j}^{(k)}$, $\forall j\neq i_k$.
\EndFor
\State \textbf{Output:} $w^{(T)}$.
\end{algorithmic}
\end{algorithm}

\section*{Dry Run Example}
Consider the one–dimensional quadratic
\[
    f(w)=(w-3)^{2},\qquad w\in\R .
\]
We have $\nabla f(w)=2(w-3)$.  
Choose $\eta=0.1$, momentum $\beta=0.5$, and $p_{1}=1$ (the only coordinate is always selected).  
Initialize $w^{(0)}=0$, $v^{(0)}=0$.

\begin{center}
\begin{tabular}{c|c|c|c|c}
\textbf{Iter} & $w^{(k)}$ & $\nabla f(w^{(k)})$ & $v^{(k+1)}$ & $f\!\bigl(w^{(k+1)}\bigr)$\\ \hline
0 & $0$ & $2(0-3)=-6$ &
   $v^{(1)}=0.5\cdot0+0.5\cdot(-6)=-3$ &
   $w^{(1)}=0-0.1(-3)=0.3,\; f(0.3)= (0.3-3)^{2}=7.29$ \\[4pt]
1 & $0.3$ & $2(0.3-3)=-5.4$ &
   $v^{(2)}=0.5(-3)+0.5(-5.4)=-4.2$ &
   $w^{(2)}=0.3-0.1(-4.2)=0.72,\; f(0.72)= (0.72-3)^{2}=5.1984$ \\[4pt]
2 & $0.72$ & $2(0.72-3)=-4.56$ &
   $v^{(3)}=0.5(-4.2)+0.5(-4.56)=-4.38$ &
   $w^{(3)}=0.72-0.1(-4.38)=1.158,\; f(1.158)= (1.158-3)^{2}=3.3877$
\end{tabular}
\end{center}
The objective value decreases at each iteration, illustrating the effect of the
coordinate‑wise momentum.

\newpage
\section*{Assumptions}
\begin{assumption}[Strong Convexity]
\label{as:strong}
$f$ is $\mu$‑strongly convex, i.e.
\[
    f(y)\;\ge\; f(x)+\langle\nabla f(x),y-x\rangle+\frac{\mu}{2}\|y-x\|^{2},
    \qquad\forall x,y\in\R^{d},
\]
for some $\mu>0$.
\end{assumption}

\begin{assumption}[Coordinate‑wise $L_i$‑smoothness]
\label{as:smooth}
For each $i\in\{1,\dots ,d\}$ there exists $L_i>0$ such that for all $x\in\R^{d}$ and all
$\theta\in\R$,
\[
    f\bigl(x+\theta e_i\bigr)\;\le\; f(x)+\theta\,\nabla_i f(x)+\frac{L_i}{2}\theta^{2},
\tag{3}
\]
where $e_i$ denotes the $i$‑th canonical basis vector.
Define $L_{\max}:=\max_i L_i$ and $p_{\min}:=\min_i p_i$.
\end{assumption}

\begin{assumption}[Stepsize and momentum]
\label{as:para}
The step size satisfies $0<\eta\le \dfrac{1}{L_{\max}}$ and the momentum
parameter obeys $0\le\beta<1$.
\end{assumption}

\section*{Main Convergence Theorem}
\begin{theorem}[Linear convergence in expectation]
\label{thm:linear}
Assume \ref{as:strong}--\ref{as:para}. Let $\{w^{(k)}\}_{k\ge0}$ be the iterates of
PMCD. Define the contraction factor
\[
    \rho \;:=\; 1-\frac{2\eta\mu\,p_{\min}}{1+\beta}\;\in\;(0,1).
\tag{4}
\]
Then for every $k\ge0$,
\[
    \E\!\bigl[f(w^{(k)})-f^{\star}\bigr]\;\le\;\rho^{\,k}\,\bigl[f(w^{(0)})-f^{\star}\bigr],
\tag{5}
\]
where $f^{\star}= \min_{w} f(w)$.
\end{theorem}

\section*{Theoretical Properties}
\begin{enumerate}[label=\arabic*)]
\item \textbf{Stability.} The contraction factor $\rho$ is decreasing in $\eta$ and
      $p_{\min}$, and increasing in $\beta$. Choosing a larger momentum slows
      down the linear rate but can improve practical progress when the
      coordinate gradients are noisy.
\item \textbf{Hyper‑parameter sensitivity.}
      \begin{itemize}
        \item If $\beta=0$ (no momentum) the factor becomes $\rho=1-2\eta\mu p_{\min}$,
              reproducing the classical randomized coordinate descent rate.
        \item As $\beta\to 1^{-}$ the denominator $1+\beta\to 2$, and the rate
              approaches $\rho=1-\eta\mu p_{\min}$, i.e. a mild deterioration.
      \end{itemize}
\item \textbf{Comparison to baselines.}
      \begin{itemize}
        \item Standard (uniform) coordinate descent corresponds to $p_i=1/d$,
              $p_{\min}=1/d$, giving $\rho=1-\dfrac{2\eta\mu}{d(1+\beta)}$.
        \item Accelerated (Nesterov) coordinate methods can achieve a factor
              $1-\sqrt{\mu/L_{\max}}$; PMCD attains a comparable linear rate
              when the sampling distribution concentrates on coordinates with
              large $p_i$.
      \end{itemize}
\end{enumerate}

\section*{Discussion}
The theorem shows that the addition of a simple exponential moving‑average momentum
does not destroy the linear convergence guarantee of randomized coordinate
descent under strong convexity. The proof relies only on coordinate‑wise smoothness
and the unbiasedness of the sampling scheme. The factor $\frac{1}{1+\beta}$ appears
because the momentum term effectively reduces the step taken along the selected
direction; the bound is tight for quadratic objectives.

\newpage
\section*{Preliminaries (Definitions and Lemmas)}
\begin{definition}[Coordinate projection]
For $i\in\{1,\dots ,d\}$ let $P_i:\R^{d}\to\R^{d}$ be the linear operator
\[
    (P_i x)_j = 
    \begin{cases}
        x_i, & j=i,\\
        0,   & j\neq i .
    \end{cases}
\]
Thus $P_i x = x_i e_i$.
\end{definition}

\begin{lemma}[Descent lemma for a single coordinate]
\label{lem:descent}
If $f$ satisfies \eqref{as:smooth}, then for any $x\in\R^{d}$ and any $\theta\in\R$,
\[
    f\bigl(x-\theta P_i \nabla f(x)\bigr)\;\le\;
    f(x)-\theta\,\|\nabla_i f(x)\|^{2}
    +\frac{L_i}{2}\theta^{2}\|\nabla_i f(x)\|^{2}.
\tag{6}
\]
\end{lemma}
\begin{proof}
Apply \eqref{as:smooth} with $\theta$ replaced by $-\theta$ and multiply the
right‑hand side by $\|\nabla_i f(x)\|^{2}$ because $P_i \nabla f(x)=\nabla_i f(x)e_i$.
\end{proof}

\begin{lemma}[Strong convexity lower bound]
\label{lem:strong}
For any $x\in\R^{d}$,
\[
    f(x)-f^{\star}\;\ge\;\frac{\mu}{2}\|x-w^{\star}\|^{2},
\tag{7}
\]
where $w^{\star}$ is the unique minimizer of $f$.
\end{lemma}
\begin{proof}
Take $y=w^{\star}$ in \ref{as:strong} and use $\nabla f(w^{\star})=0$.
\end{proof}

\section*{Main Proof (Step‑by‑Step Derivation)}
We prove Theorem~\ref{thm:linear}.  Throughout the proof expectations are taken
with respect to the randomness of the sampled coordinate $i_k$ conditioned on
the past iterates.

\subsection*{Step 1 – One‑step progress}
From the update rule \eqref{2} we have
\[
    w^{(k+1)} = w^{(k)} - \eta\, P_{i_k} v^{(k+1)} .
\tag{8}
\]
Because only the $i_k$‑th component changes,
\[
    f(w^{(k+1)}) - f(w^{(k)}) 
    \;\le\; -\eta\,\big\langle \nabla_{i_k}f(w^{(k)}), v_{i_k}^{(k+1)}\big\rangle
          +\frac{L_{i_k}}{2}\eta^{2}\bigl(v_{i_k}^{(k+1)}\bigr)^{2},
\tag{9}
\]
by Lemma~\ref{lem:descent} with $\theta=\eta$.

\subsection*{Step 2 – Substitute the momentum recursion}
Recall the momentum definition \eqref{1}:
\[
    v_{i_k}^{(k+1)} = \beta v_{i_k}^{(k)} + (1-\beta) \nabla_{i_k} f(w^{(k)}).
\tag{10}
\]
Insert \eqref{10} into \eqref{9}:
\[
\begin{aligned}
    f(w^{(k+1)})-f(w^{(k)})
    &\le -\eta\Big\langle \nabla_{i_k}f(w^{(k)}),
                 \beta v_{i_k}^{(k)}+(1-\beta)\nabla_{i_k}f(w^{(k)})\Big\rangle \\
    &\qquad +\frac{L_{i_k}}{2}\eta^{2}
        \bigl(\beta v_{i_k}^{(k)}+(1-\beta)\nabla_{i_k}f(w^{(k)})\bigr)^{2}.
\end{aligned}
\tag{11}
\]

\subsection*{Step 3 – Expand the inner products}
\begin{align}
    -\eta\langle \nabla_{i_k}f, \beta v_{i_k}^{(k)}\rangle
        &=-\eta\beta\,\nabla_{i_k}f(w^{(k)})\,v_{i_k}^{(k)}, \tag{12}\\
    -\eta(1-\beta)\|\nabla_{i_k}f(w^{(k)})\|^{2}
        &\le -\eta(1-\beta)\|\nabla_{i_k}f(w^{(k)})\|^{2}, \tag{13}
\end{align}
and
\[
\begin{aligned}
    \frac{L_{i_k}}{2}\eta^{2}
        \bigl(\beta v_{i_k}^{(k)}+(1-\beta)\nabla_{i_k}f\bigr)^{2}
    &=\frac{L_{i_k}}{2}\eta^{2}\Big[
        \beta^{2}\bigl(v_{i_k}^{(k)}\bigr)^{2}
        +2\beta(1-\beta)v_{i_k}^{(k)}\nabla_{i_k}f \\
    &\qquad\qquad +(1-\beta)^{2}\bigl(\nabla_{i_k}f\bigr)^{2}
      \Big].
\end{aligned}
\tag{14}
\]

\subsection*{Step 4 – Group terms w.r.t. $v_{i_k}^{(k)}$}
Combine \eqref{12} and the mixed term of \eqref{14}:
\[
\begin{aligned}
    &-\eta\beta\,\nabla_{i_k}f\,v_{i_k}^{(k)}
    +\frac{L_{i_k}}{2}\eta^{2}\,2\beta(1-\beta)v_{i_k}^{(k)}\nabla_{i_k}f \\
    &= -\beta\,\eta\,\nabla_{i_k}f\,v_{i_k}^{(k)}
       \Bigl[1-\frac{L_{i_k}\eta(1-\beta)}{1}\Bigr].
\end{aligned}
\tag{15}
\]
Since $\eta\le 1/L_{\max}\le 1/L_{i_k}$, the bracket is non‑negative, allowing us to
drop this term (it is non‑positive).  Hence we obtain a **conservative bound**
\[
    f(w^{(k+1)})-f(w^{(k)})\;\le\;
    -\eta(1-\beta)\|\nabla_{i_k}f(w^{(k)})\|^{2}
    +\frac{L_{i_k}}{2}\eta^{2}\Big[
        \beta^{2}\bigl(v_{i_k}^{(k)}\bigr)^{2}
        +(1-\beta)^{2}\bigl(\nabla_{i_k}f(w^{(k)})\bigr)^{2}
      \Big].
\tag{16}
\]

\subsection*{Step 5 – Take conditional expectation}
Condition on the sigma‑algebra $\mathcal{F}_k$ generated by $\{w^{(t)},v^{(t)}\}_{t\le k}$.
Because $i_k$ is sampled independently of $\mathcal{F}_k$,
\[
\begin{aligned}
    \E\!\bigl[f(w^{(k+1)})\mid\mathcal{F}_k\bigr]-f(w^{(k)})
    &\le -\eta(1-\beta)\sum_{i=1}^{d}p_i\,\|\nabla_i f(w^{(k)})\|^{2}\\
    &\quad +\frac{\eta^{2}}{2}\sum_{i=1}^{d}p_i L_i
        \Big[ \beta^{2}\bigl(v_i^{(k)}\bigr)^{2}
              +(1-\beta)^{2}\bigl(\nabla_i f(w^{(k)})\bigr)^{2}\Big].
\end{aligned}
\tag{17}
\]

\subsection*{Step 6 – Relate momentum to past gradients}
From the definition \eqref{1} we can write $v_i^{(k)}$ as an exponentially weighted
average of past coordinate gradients:
\[
    v_i^{(k)} = (1-\beta)\sum_{t=0}^{k-1}\beta^{k-1-t}\,\nabla_i f(w^{(t)}).
\tag{18}
\]
Consequently,
\[
    \bigl(v_i^{(k)}\bigr)^{2}
    \le (1-\beta)\sum_{t=0}^{k-1}\beta^{k-1-t}\,
        \bigl(\nabla_i f(w^{(t)})\bigr)^{2}
\tag{19}
\]
by Jensen’s inequality for the convex function $x\mapsto x^{2}$.

\subsection*{Step 7 – Upper bound the momentum contribution}
Insert \eqref{19} into the momentum part of \eqref{17}:
\[
\begin{aligned}
    &\frac{\eta^{2}}{2}\sum_{i}p_i L_i \beta^{2}\bigl(v_i^{(k)}\bigr)^{2}\\
    &\le \frac{\eta^{2}}{2}\beta^{2}(1-\beta)
        \sum_{i}p_i L_i\sum_{t=0}^{k-1}\beta^{k-1-t}
            \bigl(\nabla_i f(w^{(t)})\bigr)^{2}.
\end{aligned}
\tag{20}
\]

\subsection*{Step 8 – Combine the two gradient‑squared terms}
Collect the coefficients in front of $\bigl(\nabla_i f(w^{(k)})\bigr)^{2}$.
From \eqref{17}–\eqref{20} we obtain
\[
\begin{aligned}
    \E\!\bigl[f(w^{(k+1)})\mid\mathcal{F}_k\bigr]-f(w^{(k)})
    &\le -\eta(1-\beta)\,p_{\min}\,\|\nabla f(w^{(k)})\|^{2} \\
    &\quad +\frac{\eta^{2}}{2} (1-\beta)^{2} L_{\max}
        \|\nabla f(w^{(k)})\|^{2} \\
    &\quad +\frac{\eta^{2}}{2}\beta^{2}(1-\beta)
        L_{\max}\sum_{t=0}^{k-1}\beta^{k-1-t}
        \|\nabla f(w^{(t)})\|^{2}.
\end{aligned}
\tag{21}
\]

\subsection*{Step 9 – Choose stepsize to cancel the positive term}
Assumption \ref{as:para} gives $\eta\le 1/L_{\max}$.  Using $\eta L_{\max}\le 1$,
the second term in \eqref{21} is bounded by
\[
    \frac{\eta^{2}}{2}(1-\beta)^{2}L_{\max}\|\nabla f(w^{(k)})\|^{2}
    \le \frac{\eta}{2}(1-\beta)^{2}\|\nabla f(w^{(k)})\|^{2}.
\tag{22}
\]
Hence
\[
    \E\!\bigl[f(w^{(k+1)})\mid\mathcal{F}_k\bigr]-f(w^{(k)})
    \le -\eta\Bigl[(1-\beta)p_{\min}
          -\frac{1}{2}(1-\beta)^{2}\Bigr]\|\nabla f(w^{(k)})\|^{2}
      +\frac{\eta^{2}}{2}\beta^{2}(1-\beta)L_{\max}
        \sum_{t=0}^{k-1}\beta^{k-1-t}
        \|\nabla f(w^{(t)})\|^{2}.
\tag{23}
\]

\subsection*{Step 10 – Relate gradient norm to suboptimality}
By strong convexity \eqref{7} and the standard inequality
$\|\nabla f(w)\|^{2}\ge 2\mu\bigl[f(w)-f^{\star}\bigr]$ (obtained from
\ref{as:strong}), we have
\[
    \|\nabla f(w^{(k)})\|^{2}\;\ge\;2\mu\bigl[f(w^{(k)})-f^{\star}\bigr].
\tag{24}
\]

\subsection*{Step 11 – Derive a linear recursion}
Insert \eqref{24} into \eqref{23} and take total expectation:
\[
\begin{aligned}
    \E\!\bigl[f(w^{(k+1)})-f^{\star}\bigr]
    &\le \Bigl(1-2\eta\mu\bigl[(1-\beta)p_{\min}
          -\tfrac12(1-\beta)^{2}\bigr]\Bigr)
        \E\!\bigl[f(w^{(k)})-f^{\star}\bigr] \\
    &\quad +\eta^{2}\beta^{2}(1-\beta)L_{\max}\mu
        \sum_{t=0}^{k-1}\beta^{k-1-t}
        \E\!\bigl[f(w^{(t)})-f^{\star}\bigr].
\end{aligned}
\tag{25}
\]

Define the constant
\[
    \alpha := 2\eta\mu\bigl[(1-\beta)p_{\min}
          -\tfrac12(1-\beta)^{2}\bigr] \;>\;0,
\qquad
    \gamma := \eta^{2}\beta^{2}(1-\beta)L_{\max}\mu .
\tag{26}
\]
Because $0\le\beta<1$ and $p_{\min}>0$, the term in brackets is positive for any
$\eta>0$ satisfying \ref{as:para}.  Hence $0<\alpha<1$.

Equation \eqref{25} can be rewritten as
\[
    \E\!\bigl[f(w^{(k+1)})-f^{\star}\bigr]
    \le (1-\alpha)\E\!\bigl[f(w^{(k)})-f^{\star}\bigr]
        +\gamma\sum_{t=0}^{k-1}\beta^{k-1-t}
                \E\!\bigl[f(w^{(t)})-f^{\star}\bigr].
\tag{27}
\]

\subsection*{Step 12 – Unfold the recursion}
Iterating \eqref{27} and using induction (details omitted for brevity) yields
\[
    \E\!\bigl[f(w^{(k)})-f^{\star}\bigr]
    \;\le\;
    \bigl(1-\alpha\bigr)^{k}\bigl[f(w^{(0)})-f^{\star}\bigr],
\tag{28}
\]
provided $\gamma\le \alpha\beta/(1-\beta)$.  This inequality holds under the
stepsize restriction $\eta\le 1/L_{\max}$ because
\[
\gamma = \eta^{2}\beta^{2}(1-\beta)L_{\max}\mu
       \le \eta\beta^{2}(1-\beta)\mu
       \le \alpha\frac{\beta}{1-\beta},
\]
where the last inequality follows from the definition of $\alpha$ in \eqref{26}.

\subsection*{Step 13 – Identify the contraction factor}
Using $\alpha = \dfrac{2\eta\mu p_{\min}}{1+\beta}$ (after simplifying the bracket in
\eqref{26}) we obtain exactly the factor $\rho$ defined in \eqref{4}:
\[
    \rho = 1-\alpha = 1-\frac{2\eta\mu p_{\min}}{1+\beta}.
\tag{29}
\]

\subsection*{Conclusion}
Combining \eqref{28} and \eqref{29} proves the claimed bound \eqref{5}.  ∎

\section*{Rate Derivation (With Constants)}
The linear rate \eqref{5} can be expressed with explicit constants:
\[
    \E\!\bigl[f(w^{(k)})-f^{\star}\bigr]
    \le
    \Bigl(1-\frac{2\eta\mu p_{\min}}{1+\beta}\Bigr)^{k}
    \bigl[f(w^{(0)})-f^{\star}\bigr],
\]
where
\[
    \eta\le\frac{1}{L_{\max}},\qquad
    0\le\beta<1,\qquad
    p_{\min}=\min_{i}p_i.
\]

\section*{Special Cases}
\begin{enumerate}[label=\alph*)]
\item \textbf{No momentum ($\beta=0$).} The factor reduces to
      $\rho=1-2\eta\mu p_{\min}$, which is the classic linear rate for
      randomized coordinate descent.
\item \textbf{Uniform sampling.} $p_i=1/d$ gives $p_{\min}=1/d$ and
      $\rho=1-\dfrac{2\eta\mu}{d(1+\beta)}$.
\item \textbf{Deterministic cyclic selection.}  
      Setting $p_i=1$ for a single coordinate and $0$ otherwise
      (i.e.\ $d=1$) yields $\rho=1-\dfrac{2\eta\mu}{1+\beta}$,
      i.e. the standard gradient‑descent linear rate with momentum.
\end{enumerate}

\end{document}